\section{Boundary frames and centered holomorphic variations}
\label{var:section}

We establish the two analytic facts needed to vary the maximizing disc.
Throughout this section, let
\[
 f\colon\overline\D\longrightarrow X
\]
be continuous, holomorphic on $\D$, and of class $W^{1,p}$ for some
$p>2$, with the coordinate convention of Section~\ref{sec:preliminaries}. The
maximizers supplied by Proposition~\ref{ex:regularity} satisfy these
assumptions.  No extension of $f$ across $\partial\D$ is assumed.

We use two elementary consequences of $p>2$: the embedding
$W^{1,p}(\D)\subset C^0(\overline\D)$ and the fact that
$W^{1,p}(\D)$ is a Banach algebra under pointwise multiplication.
In particular, a continuous, pointwise invertible $W^{1,p}$ matrix on
$\overline\D$ has a $W^{1,p}$ inverse.

\subsection{A holomorphic frame regular at the boundary}

We first trivialize the pulled-back tangent bundle while retaining
the boundary Sobolev regularity of $f$. We will then change this
frame so that its boundary values are orthonormal for $h$.
\begin{lemma}\label{var:initial-frame}
The bundle $f^*T^{1,0}X$ admits a holomorphic frame $e$ on $\D$ which
extends to a continuous, nonsingular $W^{1,p}$ frame on $\overline\D$.
\end{lemma}

\begin{proof}
The continuous pullback bundle over $\overline\D$ is topologically
trivial.  Its coordinate transition matrices are $W^{1,p}$: they are
smooth functions of $f$ with bounded derivatives on the compact subsets
of the target charts in use.  Start with a continuous frame, approximate
its coefficients uniformly in finitely many local frames, and combine
the approximations with a smooth partition of unity on the source.
This gives a $W^{1,p}$ frame $v$ uniformly close to the original one;
sufficiently close approximation preserves its pointwise independence.

In the frame $v$, the Dolbeault operator on the open disc has the form
\[
 \bar\partial+a,
 \qquad a\in L^p(\D,\operatorname{Mat}_n(\C)),
\]
where we identify a $(0,1)$-form with its coefficient.  Extend $a$ by
zero to a disc $\D_R=\{|z|<R\}$ with $R>1$.  On a sufficiently small
disc $U\subset\D_R$ of radius $r$, solve
\begin{equation}\label{var:local-gauge}
 \bar\partial U_0=-aU_0.
\end{equation}
Here and below the same notation for a matrix equation is understood
entrywise.  For completeness, the solid Cauchy transform
\[
 T_Ug(z)=\frac1\pi\int_U\frac{g(\zeta)}{z-\zeta}\,\dd S(\zeta)
\]
satisfies $\bar\partial T_Ug=g$ weakly in $U$.  H\"older's inequality
for its kernel and the $L^p$ estimate for the planar Cauchy singular
integral~\cite[Section~4.1 of the authors' preprint]{SukhovTumanov2016} give
\begin{equation}\label{var:cauchy-estimates}
 \|T_Ug\|_{L^\infty(U)}
       \le C_p r^{1-2/p}\|g\|_{L^p(U)},
 \qquad
 \|T_Ug\|_{W^{1,p}(U)}\le C_{p,r}\|g\|_{L^p(U)}.
\end{equation}
Thus $U_0=I-T_U(aU_0)$ is solved by contraction in the sup norm when
$r$ is small.  The radius can also be chosen so that the solution stays
within distance $1/2$ of $I$.  It is therefore invertible and, by
Equation~\eqref{var:cauchy-estimates}, belongs to $W^{1,p}$.

For any two such solutions $U_i,U_j$, the transition matrix
$U_i^{-1}U_j$ satisfies $\bar\partial(U_i^{-1}U_j)=0$ weakly on the
overlap.  It is consequently holomorphic there.  These transitions
define a holomorphic vector bundle on $\D_R$.  A holomorphic vector
bundle on a disc is holomorphically trivial, by the Oka--Grauert
principle~\cite{Grauert1958,ForstnericPrezelj2000}.  Choose a global holomorphic frame of
this bundle.  Relative to its local frames $U_i$, its coefficients are
holomorphic; relative to the original trivial bundle they are therefore
$W^{1,p}$ on compact subsets of $\D_R$.  In particular the resulting
matrix $U$ is continuous, invertible, and $W^{1,p}$ on $\overline\D$,
and satisfies $\bar\partial U=-aU$ on $\D$.  The frame $e=vU$ has all
the required properties.
\end{proof}

\begin{proposition}[Boundary-unitary frame]\label{var:frame}
Let $h$ be a smooth Hermitian metric on $X$.  There exist a number
$p_0$ with $2<p_0\le p$ and a holomorphic frame
$B=(b_1,\ldots,b_n)$ of $f^*T^{1,0}X$ which is continuous,
nonsingular, and $W^{1,p_0}$ on $\overline\D$, such that
\begin{equation}\label{var:boundary-unitary}
 B(z)^*h(f(z))B(z)=I,
 \qquad z\in\partial\D.
\end{equation}
\end{proposition}

The factorization step belongs to the matrix spectral-factorization
theory of Wiener--Masani and Helson--Lowdenslager
\cite{WienerMasani1957,HelsonLowdenslager1958}. We give the weighted
Hardy-space construction and establish the boundary Sobolev regularity
required for our variations.

\begin{proof}
Choose $e$ by Lemma~\ref{var:initial-frame} and set
\[
 M(z)=e(z)^*h(f(z))e(z).
\]
This is a continuous, positive definite $W^{1,p}$ matrix on the closed
disc.  In particular, there are constants $0<m\le M_0<\infty$ such
that
\begin{equation}\label{var:metric-bounds}
 mI\le M(z)\le M_0I
 \quad\hbox{on }\overline\D.
\end{equation}

\paragraph{The Hardy-space factorization.}
Write $H^2=H^2(\partial\D,\C^n)$ for the vector Hardy space with
nonnegative Fourier modes, and put $\dd\sigma=\dd\theta/(2\pi)$.
Equip the same vector space with the equivalent Hilbert norm
\[
 \|u\|_M^2=\int_{\partial\D}u^*Mu\,\dd\sigma.
\]
Multiplication by $z$ is an isometry, and its closed image $zH^2$ has
codimension $n$.  Let $s_1,\ldots,s_n$ be an orthonormal basis of
$W=H^2\ominus_M zH^2$ and form the matrix $S=(s_1,\ldots,s_n)$.
The vectors $z^j s_a$, with $j\ge0$ and $1\le a\le n$, are an
orthonormal basis.  Indeed, iterating
$H^2=W\oplus_M zH^2$ shows that the orthogonal complement of their
span is contained in every $z^jH^2$; the intersection of these spaces
is zero.

Their orthogonality identifies all Fourier coefficients of the
$L^1$ matrix $S^*MS$, and gives
\begin{equation}\label{var:hardy-factor}
 S^*MS=I\quad\hbox{almost everywhere on }\partial\D.
\end{equation}
Equation~\eqref{var:metric-bounds} then bounds the boundary values of
$S$ in $L^\infty$.  Its holomorphic Hardy extension is bounded on
$\D$, as follows entrywise from the Poisson integral.  It is invertible
at every point $\zeta\in\D$: evaluation at $\zeta$ is a continuous
surjection from $H^2$ onto $\C^n$, whereas the evaluations of all
$z^j s_a$ lie in the range of $S(\zeta)$.  Their dense span forces
that range to be $\C^n$.

\paragraph{Regularity of the boundary factor.}
The factor constructed so far has only almost-everywhere boundary
values. We now obtain continuity and a Sobolev exponent above two,
which will also give nonsingularity at every boundary point.
We first record the fractional regularity supplied by $M\in W^{1,p}$.
If $M_\theta(e^{i\alpha})=M(e^{i(\alpha+\theta)})$, then
\begin{equation}\label{var:trace-difference}
 \|M_\theta-M\|_{L^p(\partial\D)}
       \le C|\theta|^{1-1/p},\qquad |\theta|\le\tfrac14.
\end{equation}
One can see this directly, without a trace theorem.  Set
$t=|\theta|$ and choose $\rho\in[1-2t,1-t]$ such that the angular
$L^p$ norm of $\partial_\alpha M(\rho e^{i\alpha})$ is at most
$Ct^{-1/p}\|M\|_{W^{1,p}}$.  Radial integration and H\"older's
inequality bound
$\|M|_{\partial\D}-M(\rho\,\cdot)\|_{L^p}$ by
$Ct^{1-1/p}\|M\|_{W^{1,p}}$.  At radius $\rho$, angular integration
gives the same bound for translation through $\theta$.  The triangle
inequality proves Equation~\eqref{var:trace-difference}; the argument
for Sobolev functions follows by approximation, with the continuous
boundary values identifying the trace.

Let $H^s(S^1)$ now denote the fractional Sobolev space on the circle,
defined by square summability of Fourier coefficients with weight
$1+|l|^{2s}$. This differs from the Hardy-space notation $H^2$ above.
It follows that $M|_{\partial\D}\in H^s(S^1)$ for every
\begin{equation}\label{var:s-choice}
 \frac12<s<1-\frac1p.
\end{equation}
Indeed, on a circle of finite measure the $L^p$ bound controls the
$L^2$ bound, and integration of the squared rotation differences
against $|\theta|^{-1-2s}\,\dd\theta$ is finite.  Parseval's identity
identifies this condition with
$\sum_{l\in\mathbb Z}(1+|l|^{2s})|M_l|^2<\infty$.

Let $\Pi$ denote the unweighted orthogonal Hardy projection.  The
Toeplitz operator
\[
 T_Mv=\Pi(Mv),\qquad T_M\colon H^2\longrightarrow H^2,
\]
is bounded, self-adjoint, and satisfies
$\langle T_Mv,v\rangle\ge m\|v\|_2^2$.  Consequently it has a
bounded inverse: its range is closed by the lower bound and dense
because its orthogonal complement is the kernel of its adjoint.
For a column $v$ of $S$, membership in $W$ says precisely that
$T_Mv$ is a constant vector.  Rotating this identity and subtracting
it from the original one yields
\[
 T_M(v_\theta-v)=\Pi\bigl((M-M_\theta)v_\theta\bigr).
\]
The boundedness of $v$ and $T_M^{-1}$ gives
\begin{equation}\label{var:factor-difference}
 \|v_\theta-v\|_2
 \le C\|(M_\theta-M)v_\theta\|_2
 \le C\|M_\theta-M\|_2.
\end{equation}
The same rotation characterization therefore gives $v\in H^s$.

Write $v(z)=\sum_{l\ge0}v_lz^l$.  Since $s>1/2$, Cauchy--Schwarz
shows that this series converges absolutely and uniformly on
$\overline\D$.  Thus $v$ is continuous there.  More precisely, for
$1/2<r<1$, Parseval's identity and Cauchy--Schwarz give
\begin{align}
 \|v'(r\,\cdot)\|_2^2
 &\le C(1-r)^{2s-2}
       \sum_{l\ge1}l^{2s}|v_l|^2,
       \label{var:factor-l2}\\
 \|v'(r\,\cdot)\|_\infty
 &\le C(1-r)^{s-3/2}
       \left(\sum_{l\ge1}l^{2s}|v_l|^2\right)^{1/2}.
       \label{var:factor-sup}
\end{align}
For any $p_0>2$, interpolation of these two estimates bounds
$\|v'(r\,\cdot)\|_{p_0}^{p_0}$ by a constant times
\[
 (1-r)^{(s-3/2)(p_0-2)+2s-2}
       =(1-r)^{p_0(s-3/2)+1}.
\]
This is integrable in $r$ provided
\begin{equation}\label{var:p0-choice}
 2<p_0<\frac{2}{3/2-s}.
\end{equation}
The interval is nonempty because $s>1/2$, and we choose $p_0$ also
at most $p$.  It follows that $S\in W^{1,p_0}(\D)$.

By continuity, Equation~\eqref{var:hardy-factor} now holds at every
boundary point.  It implies that $S$ is invertible there as well as
in the interior.  Compactness gives uniform nonsingularity on
$\overline\D$.  Finally $B=eS$ is a continuous, nonsingular,
$W^{1,p_0}$ holomorphic frame and satisfies
Equation~\eqref{var:boundary-unitary}.  Notice that uniform
nonsingularity was obtained after boundary regularity; it was not
assumed in the Hardy-space argument.
\end{proof}

\subsection{Integrating the infinitesimal variations}

The boundary-unitary frame provides sections $zb_j$ vanishing at the
center. To use maximality, these sections must be derivatives of
admissible holomorphic discs. Sprays of analytic discs fixing an
exceptional set are standard in disc-envelope theory
\cite[Definition~2.1 and Lemma~2.2]{DrinovecForstneric2012}.
The following construction records the prescribed derivatives and
Sobolev parameter dependence required by the Hessian calculation.

\begin{proposition}[Centered holomorphic variations]\label{var:spray}
Let $B$ and $p_0>2$ be as in Proposition~\ref{var:frame}, and set
$x=f(0)$.  For some $\eta>0$, there is a map
\[
 F\colon\{t\in\C^n:|t|<\eta\}\times\overline\D\longrightarrow X
\]
which is continuous on its domain and jointly holomorphic on
$\{|t|<\eta\}\times\D$, and satisfies
\begin{equation}\label{var:spray-identities}
 F(0,z)=f(z),\qquad F(t,0)=x,\qquad
 \frac{\partial F}{\partial t_j}(0,z)=z b_j(z).
\end{equation}
In fixed local target coordinates covering the graph of $f$, the
maps $t\mapsto F(t,\cdot)$ are holomorphic with values in
$W^{1,p_0}$ on the corresponding source patches.  In particular
every $F(t,\cdot)$ is a finite-area holomorphic disc centered at $x$.
\end{proposition}

The proof first gives the graph of $f$ fiberwise holomorphic
coordinates with Sobolev dependence on the source. The resulting
coordinate maps need not be holomorphic in the source; we correct
that defect by a nonlinear $\bar\partial$ equation. Subtracting the
central value in its right inverse preserves the prescribed center.

\begin{proof}
\textbf{Fiber coordinates.}
Choose a frame $e$ from Lemma~\ref{var:initial-frame} and write
$B=eS$. The matrix $S=e^{-1}B$ is holomorphic on $\D$ and belongs
to $W^{1,p_0}$ on $\overline\D$. We first construct fiber coordinates near the graph of
$f$.  Choose a finite cover of $\overline\D$ by relatively open
sets $U_i$ and target charts
$\phi_i\colon V_i\longrightarrow\C^n$ so that $f(\overline U_i)$
is a compact subset of $V_i$, after shrinking the source sets as
necessary.  Let $\rho_i$ be a smooth partition of unity on the
closed disc subordinate to these sets.  In the $i$th chart put
\[
 A_i(z)=D\phi_i(f(z))e(z),\qquad
 C_i(z,y)=A_i(z)^{-1}\bigl(\phi_i(y)-\phi_i(f(z))\bigr).
\]
For $z\in U_i$, the matrix $A_i$ is holomorphic in the interior,
continuous and invertible up to the boundary, and of class $W^{1,p_0}$.
The map $C_i$ is holomorphic in $y$ and, for fixed admissible $y$,
holomorphic in the interior source variable $z$.  On a neighborhood
of the graph of $f$ define
\begin{equation}\label{var:forward-coordinates}
 \Phi(z,y)=\sum_i\rho_i(z)C_i(z,y).
\end{equation}
Each term is used only where its target chart is defined; the compact
support of $\rho_i$ inside the corresponding source patch makes
extension by zero harmless.  These neighborhoods can be chosen
uniformly about the graph.  More precisely, on sufficiently small
source patches all indices whose supports meet the patch have their
charts defined on one common target neighborhood of its image under
$f$.  This follows from finiteness of the cover and subordination of
the supports.  By construction,
\begin{equation}\label{var:coordinate-normalization}
 \Phi(z,f(z))=0,
 \qquad D_y\Phi(z,f(z))\,e(z)=I.
\end{equation}

\paragraph{Uniform inverses and their Sobolev dependence.}
Work on one of these smaller source patches, with fixed target chart
$\psi$, and put $c(z)=\psi(f(z))$ and
$e_\psi(z)=D\psi(f(z))e(z)$.  Normalize the forward map by
\[
 K_z(v)=\Phi\bigl(z,\psi^{-1}(c(z)+e_\psi(z)v)\bigr).
\]
Then $K_z(0)=0$ and $D_vK_z(0)=I$.  Continuity of all fiber
derivatives and compactness allow a common $R>0$ such that the maps
are defined on $|v|\le R$ and
\[
 \|D_vK_z(v)-I\|\le\tfrac14.
\]
For $|\xi|<R/2$, the equation
\[
 v=\xi-\bigl(K_z(v)-v\bigr)
\]
is a contraction of the closed ball of radius $R$ into itself.
Its solution $V(z,\xi)$ is continuous in $(z,\xi)$ and holomorphic
in $\xi$, as follows from its uniformly convergent contraction
iterations.  Define $E$ by
$\psi(E(z,\xi))=c(z)+e_\psi(z)V(z,\xi)$.  The local inverses
agree near $\xi=0$ by uniqueness and then throughout their common
fiber domain by holomorphy.  Taking the minimum of the finitely many
radii gives a common fiber neighborhood, with
\[
 \Phi(z,E(z,\xi))=\xi,
\]
and
\begin{equation}\label{var:inverse-normalization}
 E(z,0)=f(z),\qquad E_\xi(z,0)=e(z).
\end{equation}

We next obtain a Sobolev estimate uniform in the fiber parameter.
The chart formula for $K_z$ depends smoothly on the finite coefficient list
\[
 a(z)=\bigl(\rho_i(z),A_i(z)^{-1},\phi_i(f(z)),
                 c(z),e_\psi(z)\bigr)_i.
\]
All these coefficients are $W^{1,p_0}$ with compact range.  Fix a
smaller closed fiber polydisc inside $|\xi|<R/2$.  The strict bounds
in the contraction argument persist when the coefficient list varies
in a sufficiently small neighborhood of its compact range.  The
finite-dimensional implicit function theorem, with the coefficients
regarded as real parameters, therefore makes $V$ a smooth function of
that list and $\xi$, holomorphic in $\xi$.  After shrinking this
coefficient neighborhood, its first coefficient derivatives are
uniformly bounded on the fixed fiber polydisc.  The same is true for
the coordinate inverse
$E_\psi(z,\xi):=\psi(E(z,\xi))=c(z)+e_\psi(z)V(z,\xi)$.
The Sobolev chain rule now shows that $E_\psi(\cdot,\xi)$ is
$W^{1,p_0}$, with its weak source derivatives bounded by a single
$L^{p_0}$ function independent of $\xi$.  In particular, for some
$b\in L^{p_0}$ on the source patch,
\begin{equation}\label{var:inverse-sobolev-bound}
 |E_\psi(z,\xi)|\le C,
 \qquad |\nabla_zE_\psi(z,\xi)|\le C b(z)
 \quad\hbox{for almost every }z,
\end{equation}
uniformly for $\xi$ in that smaller polydisc.  One may take $b$ to
be a constant plus the sum of the absolute values of the first weak
derivatives of the finite coefficient list.  Derivatives with respect
to $\xi$ satisfy the same type of bound after another fixed shrink
of the polydisc, by Cauchy's formula.

\paragraph{Holomorphic substitution.}
These estimates also justify holomorphy of the substitution operator
\begin{equation}\label{var:inverse-substitution}
 u\longmapsto E_\psi(\cdot,u(\cdot))
\end{equation}
from a neighborhood of zero in $W^{1,p_0}(\D,\C^n)$ into the local
$W^{1,p_0}$ space.  Indeed, expand in the fiber variable:
\[
 E_\psi(z,\xi)=\sum_{\alpha\in\mathbb N^n}E_\alpha(z)\xi^\alpha.
\]
Cauchy's formula on a fixed fiber polydisc, with weak $z$
differentiation under its contour integral, gives
\[
 \|E_\alpha\|_\infty+\|\nabla E_\alpha\|_{p_0}
       \le Cr_1^{-|\alpha|}
\]
for some $r_1>0$, by Equation~\eqref{var:inverse-sobolev-bound}.
The algebra norm $\|a\|_\infty+\|\nabla a\|_{p_0}$ is submultiplicative.
Thus $\sum_\alpha E_\alpha u^\alpha$ converges normally in the
local $W^{1,p_0}$ space when the global $W^{1,p_0}$ norm of $u$ is small;
the number of multi-indices of a fixed degree grows polynomially.
Each term is a continuous
homogeneous polynomial in $u$, proving the assertion.  This argument
only concerns fixed local target coordinates; no linear structure on
$X$ is being presumed.

\paragraph{The nonlinear holomorphic-map equation.}
The fiber coordinates now have the necessary Sobolev regularity.
It remains to correct their antiholomorphic derivative in the source
while keeping the center and the prescribed first-order directions.
For $u\in W^{1,p_0}(\D,\C^n)$ close to zero, the Sobolev chain rule
in a target chart gives
\[
 \bar\partial\bigl(E(z,u(z))\bigr)
       =E_\xi(z,u(z))\bar\partial u
          +\bar\partial_zE(z,u(z)).
\]
There is no $\bar\partial\overline u$ term, since $E$ is holomorphic
in $\xi$.  Thus the map is weakly holomorphic precisely when
\begin{equation}\label{var:nonlinear-equation}
 \bar\partial u+Q(\cdot,u)=0,
 \qquad
 Q(z,\xi)=E_\xi(z,\xi)^{-1}\bar\partial_zE(z,\xi).
\end{equation}
The formula is independent of the chosen target chart, because the
derivative of a holomorphic chart change multiplies both terms.

Differentiating $\Phi(z,E(z,\xi))=\xi$ with $\xi$ fixed and using
$E_\xi=(D_y\Phi)^{-1}$ gives
\begin{equation}\label{var:q-forward}
 Q(z,\xi)=-\bar\partial_z\Phi(z,E(z,\xi))
       =-\sum_i(\bar\partial\rho_i)(z)C_i(z,E(z,\xi)).
\end{equation}
For the second equality we used holomorphy in $z$ of $C_i$ with $y$
fixed.  In particular $Q$ is holomorphic in $\xi$ and jointly
continuous and bounded on the closed disc times a smaller closed
fiber polydisc.  Cauchy's formula gives coefficients satisfying
$\|Q_\alpha\|_\infty\le Cr_2^{-|\alpha|}$ for some $r_2>0$.
The normally convergent series $\sum_\alpha Q_\alpha u^\alpha$
therefore shows that
\begin{equation}\label{var:q-analytic}
 \mathcal Q\colon W^{1,p_0}(\D,\C^n)\longrightarrow L^{p_0}(\D,\C^n),
 \qquad \mathcal Q(u)=Q(\cdot,u(\cdot)),
\end{equation}
is holomorphic near zero, using the embedding of $W^{1,p_0}$ into
the sup norm.

The constant and linear terms vanish.  Indeed,
$C_i(z,f(z))=0$ and
$D_yC_i(z,f(z))e(z)=I$.  Equations~\eqref{var:inverse-normalization}
and~\eqref{var:q-forward}, together with
$\sum_i\bar\partial\rho_i=0$, therefore give
\begin{equation}\label{var:q-vanishing}
 Q(z,0)=0,\qquad Q_\xi(z,0)=0.
\end{equation}
Equivalently, $\mathcal Q(0)=0$ and $D\mathcal Q(0)=0$ as maps
between the indicated Banach spaces.

\paragraph{The implicit function theorem with the center fixed.}
On the full disc, let $T\colon L^{p_0}\to W^{1,p_0}$ be the solid Cauchy
right inverse from Equation~\eqref{var:cauchy-estimates}.  Since
point evaluation is bounded on $W^{1,p_0}$, the operator
\[
 Hg=Tg-(Tg)(0)
\]
is bounded and satisfies
\begin{equation}\label{var:centered-inverse}
 \bar\partial Hg=g,\qquad (Hg)(0)=0.
\end{equation}
Seek a solution of Equation~\eqref{var:nonlinear-equation} in the
form
\begin{equation}\label{var:u-ansatz}
 u(t,z)=zS(z)t+Hg(t)(z),\qquad g(t)\in L^{p_0}(\D,\C^n).
\end{equation}
Since $S$ is holomorphic and $W^{1,p_0}$, the equation becomes
\[
 G(g,t):=g+\mathcal Q\bigl(zSt+Hg\bigr)=0.
\]
This is a holomorphic equation with values in $L^{p_0}$, and
Equation~\eqref{var:q-vanishing} gives
\[
 G(0,0)=0,\qquad D_gG(0,0)=I.
\]
The complex Banach implicit function
theorem~\cite[Theorem~2.3(ii), (d) of the author's preprint]{Glockner2006}
supplies a holomorphic
solution $g(t)$ for $t$ in a neighborhood of zero.  Differentiating
the equation at zero gives
\[
 g(0)=0,\qquad D_tg(0)=0.
\]
Consequently $u(t)$ is holomorphic with values in $W^{1,p_0}$, is small
in the sup norm for $t$ sufficiently small, and satisfies
\[
 u(t,0)=0,\qquad u(0,z)=0,\qquad
 \frac{\partial u}{\partial t_j}(0,z)=zS(z)e_j,
\]
where $e_j$ here denotes the $j$th standard vector of $\C^n$.

Set $F(t,z)=E(z,u(t,z))$.  The identities in
Equation~\eqref{var:spray-identities} follow from
Equation~\eqref{var:inverse-normalization} and $eS=B$.
Equation~\eqref{var:inverse-substitution} gives its holomorphic
parameter dependence in local $W^{1,p_0}$ coordinates.  In particular
it is jointly continuous up to the source boundary.  For each fixed
$t$, Equation~\eqref{var:nonlinear-equation} makes its coordinate
functions weakly holomorphic on the open disc.  They are therefore
ordinary holomorphic functions.  For each fixed $z$, bounded point
evaluation on $W^{1,p_0}$ and holomorphy of $E$ in its fiber variable
give holomorphy in $t$.  Local separate holomorphy now gives joint
holomorphy on $\{|t|<\eta\}\times\D$.

Finally $F(t,\cdot)$ is $W^{1,p_0}$ with $p_0>2$ and has compact image,
so its area with respect to any smooth Hermitian metric on $X$ is
finite.  Thus all parameters in a sufficiently small complex
neighborhood give admissible centered discs.  The construction imposes
no condition on their boundary maps and requires no extension across
$\partial\D$.
\end{proof}

\begin{remark}\label{var:center-derivative}
Although derivative evaluation at a point is not bounded on the full
space $W^{1,p_0}$, it is bounded on its holomorphic subspace on each
smaller disc, by the Cauchy formula.  In particular
$t\mapsto \partial_zF(t,0)$ is holomorphic.  Hence the central
derivative functional can be differentiated along
Proposition~\ref{var:spray} without additional boundary regularity.
\end{remark}
